\section{Related work}
\paragraph{Image-conditioned 3D generation.}
TRELLIS.2 learns a compact native representation of geometry and material
attributes using sparse voxels~\cite{xiang2025trellis2}.
Pixal3D lifts image features into a back-projected conditioning volume and
extends this mechanism to multiple views~\cite{li2026pixal3d}.
Our implementation continues the supplied multi-view Pixal3D checkpoints
and adapts their source-conditioning path. The separately implemented
multi-view prior and checkpoint lineage are specified in our protocol;
we do not identify every implementation detail with the original paper branch.

\paragraph{Image-conditioned editing.}
ShapeUP learns supervised latent translation from a source 3D shape and an
edited image to an edited 3D shape~\cite{gat2026shapeup}.
Our evaluated input instead consists of an edited target image and undeformed
source images with calibrated cameras. This input distinction motivates our
study of appearance evidence under conflicting geometric states. It does not
establish a performance advantage over source-shape-conditioned editing;
no such quantitative comparison has yet been executed in this study.

\paragraph{Representations for correspondence.}
Zhang et al. examine complementary correspondence properties of diffusion
and DINO features~\cite{zhang2023twofeatures}.
VGGT-$\Omega$ learns geometry-aware representations through reconstruction
training~\cite{wang2026vggtomega}. RoMa-$\Omega$ distinguishes zero-shot
feature matching from learned matching on feed-forward 3D
representations~\cite{nordstrom2026romaomega}.
These works motivate separating raw descriptor quality, learned matching,
and candidate aggregation. Our experiments additionally ask whether their
correspondence evidence improves generated geometry and source appearance
when delivered through the pretrained projection-conditioning path.
